\documentclass[11pt]{article}
% Shared print layout. The web's provisional design and MathJax are independent.
\usepackage[utf8]{inputenc}
\usepackage[T1]{fontenc}
\usepackage[spanish,es-nodecimaldot]{babel}
\usepackage{lmodern}
\usepackage{amsmath,amssymb}
\usepackage{graphicx}
\usepackage{geometry}
\geometry{a4paper,left=20mm,right=20mm,top=25mm,bottom=22mm,headheight=16pt,headsep=8mm}
\usepackage{fancyhdr,lastpage}
\usepackage{enumitem}
\usepackage{xcolor}
\usepackage{pgfplots}
\pgfplotsset{compat=1.18}
\usepgfplotslibrary{fillbetween}
\usepackage[hidelinks,pdfencoding=auto]{hyperref}
\hypersetup{pdflang=es-ES,pdfauthor={},pdfsubject={Matemáticas Aplicadas a las Ciencias Sociales II}}
\pagestyle{fancy}
\fancyhf{}
\fancyhead[L]{\small Matemáticas Aplicadas a las Ciencias Sociales II}
\fancyhead[R]{\small 2.\textsuperscript{o} Bachillerato}
\fancyfoot[L]{\footnotesize Distribuciones II | Edición revisada 1.0}
\fancyfoot[R]{\footnotesize Página \thepage\ de \pageref{LastPage}}
\renewcommand{\headrulewidth}{0.4pt}
\setlength{\parindent}{0pt}
\setlength{\parskip}{0.5em}
\setlist{itemsep=0.3em,topsep=0.3em}
\newcommand{\apartado}[1]{\subsection*{#1}}
\pgfplotsset{
  revision/.style={
    width=0.94\linewidth,height=7.5cm,
    axis lines=middle,grid=major,grid style={gray!18},
    tick label style={font=\small},label style={font=\small},
    xlabel={$x$},ylabel={$y$},samples=160,
    legend style={at={(0.5,-0.16)},anchor=north,draw=none,font=\small},
    clip=true
  }
}

\hypersetup{pdftitle={Distribuciones II - Examen revisado}}
\begin{document}
\begin{minipage}{0.17\linewidth}\includegraphics[width=0.8\linewidth]{logo-al-andalus.png}\end{minipage}\hfill
\begin{minipage}{0.79\linewidth}{\Large\bfseries Distribuciones II}\\[0.4em]
Duración: \textbf{55 minutos}. Puntuación: \textbf{10 puntos}.\end{minipage}

Apellidos y nombre: \hrulefill

Fecha: \hrulefill

Responde razonadamente y numera tus páginas. Se permite calculadora. Cada ejercicio incluye \textbf{0,5 puntos} para valorar la expresión de procesos y métodos; no se añaden a su puntuación máxima. Se usa $N(\mu,\sigma)$: el segundo parámetro es la desviación típica, conforme a los cálculos originales.
\section*{Ejercicio 1 \hfill 4 puntos}
En una universidad, uno de cada cinco estudiantes utiliza bicicleta. Se elige a ocho al azar y $X$ cuenta cuántos de ellos la utilizan.
\begin{enumerate}[label=(\alph*)]
\item (1,5 puntos) Identifica la distribución y calcula $P(X=0)$.
\item (1,5 puntos) Calcula la probabilidad de que al menos seis usen bicicleta.
\item (1 punto) Calcula la probabilidad de que al menos uno use bicicleta.
\end{enumerate}
\section*{Ejercicio 2 \hfill 4 puntos}
El rendimiento anual de los empleados de una empresa financiera sigue una distribución normal $N(70,15)$.
\begin{enumerate}[label=(\alph*)]
\item (1,5 puntos) Calcula el \textbf{porcentaje} con puntuación entre 60 y 90.
\item (1,5 puntos) Halla la probabilidad de obtener menos de 50 puntos.
\item (1 punto) Determina la puntuación mínima de bonificación si solo el 18\% superior la recibe.
\end{enumerate}
\section*{Ejercicio 3 \hfill 2 puntos}
Una compañía lanza 100 productos distintos. Cada uno tiene probabilidad del 60\% de ser exitoso. Para obtener beneficios deben ser exitosos al menos 70. Calcula la probabilidad de lograrlo.
\end{document}
